\section{Kết luận và hướng phát triển}

\subsection{Mức độ hoàn thành yêu cầu bài tập lớn}
Dự án đã hoàn thành các yêu cầu cốt lõi của đề bài dữ liệu văn bản:
\begin{itemize}[nosep]
    \item Thực hiện EDA và trực quan hóa phân phối dữ liệu.
    \item Xây dựng pipeline tiền xử lý văn bản có khả năng tái sử dụng.
    \item Triển khai tối thiểu hai hướng trích xuất đặc trưng (truyền thống và hiện đại).
    \item Lưu đặc trưng ra file \texttt{.npz/.npy} để tái dùng cho huấn luyện.
    \item Huấn luyện và so sánh nhiều mô hình phân loại với bộ chỉ số đầy đủ.
\end{itemize}

\subsection{Kết quả chính}
Các kết luận định lượng quan trọng:
\begin{enumerate}[nosep]
    \item Mô hình tốt nhất: \textbf{TF-IDF Bigram + LinearSVC (C=0.5)} với F1-weighted = 0.8631.
    \item Tuning tham số giúp cải thiện hiệu năng nhỏ nhưng ổn định (SVM +0.0032; LR +0.0024).
    \item Thực nghiệm cho thấy các embedding hiện đại (BERT/Word2Vec dạng feature extractor) chưa vượt qua TF-IDF bigram trong bài toán này.
    \item Sai số tập trung ở các lớp có ngữ nghĩa gần nhau, tiêu biểu \texttt{Credit card}.
\end{enumerate}

\subsection{Đóng góp kỹ thuật của hệ thống}
Bên cạnh kết quả mô hình, hệ thống có các đóng góp thực tiễn:
\begin{itemize}[nosep]
    \item Thiết kế module rõ ràng, tách được preprocessing/feature/training.
    \item Lưu đầy đủ artifact phục vụ tái lập thí nghiệm.
    \item Có sẵn unit tests cho các module trọng yếu.
    \item Notebook đáp ứng kịch bản chạy end-to-end trên nhiều nền tảng điện toán khác nhau.
\end{itemize}

\subsection{Hạn chế hiện tại}
\begin{itemize}[nosep]
    \item BERT mới dùng ở chế độ embedding tĩnh, chưa fine-tune theo miền CFPB.
    \item Chưa thực hiện calibration hoặc threshold optimization theo lớp.
\end{itemize}

\subsection{Hướng phát triển tiếp theo}
Các hướng mở rộng khả thi:
\begin{enumerate}[nosep]
    \item Fine-tune transformer trực tiếp trên bài toán phân loại 7 lớp.
    \item Thử nghiệm hybrid feature (TF-IDF + contextual embedding) và stacking.
    \item Mở rộng đánh giá với k-fold stratified và kiểm định thống kê chênh lệch mô hình.
    \item Thiết lập quy trình CI để tự động kiểm thử, tái sinh artifact và kiểm tra tính nhất quán dữ liệu.
\end{enumerate}
